% ============================================================================
% Fichier  : partie4/P04_C16_audit_securite_infrastructure.tex
% Titre    : Auditer la Sécurité et l'Infrastructure
% Auteur   : MINKA MI NGUIDJOÏ Thierry Emmanuel
% Version  : 1.0 -- octobre 2026
% ============================================================================
% Sources : loi n° 2010/012 (texte français consulté) : art. 13 et 32 (audit
% légal, rappel du chapitre 12), art. 69 (accès sans droit) et 70 (attaque
% par saturation) cités pour l'encadrement des tests ; ISO/IEC 27001:2022
% (clauses 4 à 10, § 9.2 audit interne, déclaration d'applicabilité),
% ISO/IEC 27007:2020 ; ISO/IEC 27002:2022 (numéros 8.8, 8.9, 8.15, 8.20,
% 8.22, 8.24 à vérifier dans l'édition en vigueur) ; guide du CHAI, fiche
% 3.2 « audit de sécurité » ; TP « infrastructure » du cours Audit_SI de
% l'auteur (VLAN, matrice de flux, comptes par défaut, protocoles en clair),
% mis en cohérence. Aucun scan réel : les exports du ACSI-LAB-08 sont figés.
% ============================================================================

\chapter{Auditer la Sécurité et l'Infrastructure}
\label{chap:audit-securite}
\chaptermark{Auditer la sécurité}

\epigraph{%
    « Une politique de sécurité dit ce qui devrait être fermé.\\
    Une configuration dit ce qui l'est. »%
}{Adage des auditeurs techniques}

\niveauChapitre{Intermédiaire}{%
    Maîtrise du Chapitre~\ref{chap:referentiels} (ISO/IEC~27001 et 27002)
    et du Chapitre~\ref{chap:droit} (loi n°~2010/012)%
}

\begin{boxobjectifs}
\begin{itemize}
    \item Distinguer trois audits de sécurité~: l'audit légal, l'audit d'un
          SMSI et l'audit technique.
    \item Auditer un SMSI selon l'ISO/IEC~27007, de l'appréciation des
          risques à la preuve de fonctionnement des mesures.
    \item Décomposer l'audit d'une infrastructure~: segmentation,
          durcissement, protocoles, vulnérabilités, journalisation.
    \item Comparer les flux autorisés et les flux observés.
    \item Encadrer un test technique~: autorisation écrite, périmètre,
          méthodes, preuves.
\end{itemize}
\end{boxobjectifs}

\bigskip

% ============================================================================
\section*{Question directrice : ce qui est écrit est-il configuré~?}
% ============================================================================

LiZbiZ exploite quatre centres de données, à Douala, Yaoundé, Garoua et
Bafoussam. Sa politique de sécurité interdit les protocoles
d'administration en clair, impose le changement des mots de passe par
défaut et sépare les postes des utilisateurs des serveurs de base de
données. Le domaine~D a montré que les accès étaient mal tenus
(Chapitre~\ref{chap:decomposition}). AUDIX CONSEIL doit maintenant dire si
l'infrastructure est configurée comme la politique l'écrit.

À Douala, la BANQUE SANAGA prépare à la fois l'audit de sécurité annuel de
l'ANTIC et la certification ISO/IEC~27001 de son centre de données.

Ce chapitre distingue les trois audits de sécurité
(section~\ref{sec:c16-trois}), présente l'audit d'un SMSI
(section~\ref{sec:c16-smsi}), décompose l'audit de l'infrastructure
(section~\ref{sec:c16-infra}) et encadre les tests techniques
(section~\ref{sec:c16-tests}). Le ACSI-LAB-08 applique la méthode aux exports
figés des quatre centres de LiZbiZ.

\separateur

% ============================================================================
\section{Trois Audits de Sécurité}
\label{sec:c16-trois}
% ============================================================================

\begin{tableaucours}{L{2.4cm}L{3.2cm}L{3.4cm}L{3.2cm}}{%
    \tableauHeader{Audit} &
    \tableauHeader{Fondement} &
    \tableauHeader{Objet} &
    \tableauHeader{Destinataire}
}
Audit de sécurité légal & Loi n°~2010/012, art.~13 et~32 & Politique,
mesures, procédures et moyens, au regard des risques & Ministre chargé des
Télécommunications~; rapport confidentiel \\
Audit d'un SMSI & ISO/IEC~27001 (§~9.2) et 27007~; certification par un
organisme accrédité & Le système de management~: risques, mesures,
amélioration & Direction~; organisme de certification \\
Audit technique & Lettre de mission et note d'autorisation & Configurations,
vulnérabilités, flux, comptes & RSSI, DSI \\
\end{tableaucours}
\vspace{-0.8em}
\begin{center}{\small\itshape Tableau~16.1. Trois audits de sécurité}\end{center}

Les trois se complètent. L'audit d'un SMSI vérifie que l'organisation
\emph{gère} sa sécurité~; l'audit technique vérifie que les mesures sont
\emph{effectivement en place}~; l'audit légal, défini largement par
l'article~4 de la loi (Chapitre~\ref{chap:droit}), couvre les deux. Le
guide du CHAI rappelle en outre que le dispositif de sécurité doit être
justifié par les actifs qu'il protège~: l'auditeur examine aussi
l'inventaire des actifs et l'appréciation des risques qui l'ont fondé.

% ============================================================================
\section{Auditer un SMSI}
\label{sec:c16-smsi}
% ============================================================================

\subsection{La chaîne du SMSI}

L'ISO/IEC~27001 organise ses exigences en clauses, de la clause~4
(contexte) à la clause~10 (amélioration)~: contexte et périmètre,
leadership, planification (appréciation et traitement des risques),
support, fonctionnement, évaluation des performances (dont l'audit
interne et la revue de direction), amélioration~\cite{iso27001_2022}.
L'ISO/IEC~27007 guide l'auditeur clause par clause~\cite{iso27007_2020}.

Pour l'auditeur, le SMSI est une chaîne, et chaque maillon laisse une
pièce.

\begin{tableaucours}{L{3.6cm}L{4.4cm}L{4.2cm}}{%
    \tableauHeader{Maillon} &
    \tableauHeader{Pièce} &
    \tableauHeader{Question d'audit}
}
Périmètre & Déclaration de périmètre & Les systèmes critiques y
figurent-ils~? \\
Appréciation des risques & Registre des risques de sécurité & Est-elle
datée, revue, cohérente avec les actifs~? \\
Traitement & Plan de traitement, déclaration d'applicabilité & Chaque
exclusion de mesure est-elle justifiée~? \\
Mesures & Procédures, configurations, journaux & Chaque mesure retenue
fonctionne-t-elle~? \\
Évaluation & Rapports d'audit interne, revue de direction & Les écarts
sont-ils suivis jusqu'à leur clôture~? \\
\end{tableaucours}
\vspace{-0.8em}
\begin{center}{\small\itshape Tableau~16.2. La chaîne du SMSI et ses
preuves}\end{center}

\begin{boxinfo}
\textbf{Une mesure déclarée n'est pas une mesure en place.} La déclaration
d'applicabilité dit qu'une mesure est retenue~; elle ne dit pas qu'elle
fonctionne. Sans preuve de fonctionnement, l'auditeur cote \NEK{}, et le
taux d'incertitude du SMSI mesure la distance entre le papier et la
réalité.
\end{boxinfo}

% ============================================================================
\section{Auditer l'Infrastructure}
\label{sec:c16-infra}
% ============================================================================

\subsection{Cinq sous-domaines}

\begin{tableaucours}{cL{3.0cm}L{5.6cm}L{2.2cm}}{%
    \tableauHeader{Code} &
    \tableauHeader{Sous-domaine} &
    \tableauHeader{Contrôle clé} &
    \tableauHeader{ISO/IEC 27002}
}
IN.1 & Architecture et segmentation & Seuls les flux autorisés par la
matrice de flux traversent les zones & 8.20, 8.22 \\
IN.2 & Durcissement & Aucun compte par défaut actif avec son mot de passe
d'origine & 8.9 \\
IN.3 & Protocoles & Aucun protocole d'administration ou de transfert en
clair & 8.20, 8.24 \\
IN.4 & Vulnérabilités et correctifs & Les correctifs critiques sont
appliqués dans le délai de la politique & 8.8 \\
IN.5 & Journalisation et surveillance & Les équipements envoient leurs
journaux à un point central surveillé & 8.15 \\
\end{tableaucours}
\vspace{-0.8em}
\begin{center}{\small\itshape Tableau~16.3. Décomposition de l'audit de
l'infrastructure (numéros de l'édition 2022, à vérifier)}\end{center}

IN.5 rejoint l'article~29 de la loi n°~2010/012 (surveillance des accès,
conservation des journaux, Chapitre~\ref{chap:droit}) et le sous-domaine
D.5 du contrôle d'accès~: l'auditeur réutilise les mêmes preuves.

\subsection{Flux autorisés, flux observés}

La segmentation se vérifie en comparant deux matrices~: celle des flux que
la politique autorise entre zones, et celle des flux que les règles des
pare-feu laissent réellement passer. L'écart entre les deux est le
constat.

\begin{tableaucours}{L{3.0cm}cccc}{%
    \tableauHeader{De \textbackslash{} vers} &
    \tableauHeader{DMZ} & \tableauHeader{Applicatif} &
    \tableauHeader{Utilisateurs} & \tableauHeader{Administration}
}
Internet & 443 & -- & -- & -- \\
DMZ & -- & 8443 & \textbf{445 (observé)} & -- \\
Utilisateurs & 443 & 8443 \textbf{+ 1433 (observé)} & -- & -- \\
Administration & 22 & 22 & 22 & -- \\
\end{tableaucours}
\vspace{-0.8em}
\begin{center}{\small\itshape Tableau~16.4. Matrice de flux d'un centre de
données~: ports autorisés, et en gras les flux observés non
autorisés}\end{center}

Dans cet exemple, deux flux non autorisés sont ouverts~: de la DMZ vers les
postes (un serveur web compromis pourrait atteindre les utilisateurs), et
des postes vers la base de données (port~1433~: un utilisateur peut
contourner l'application et attaquer la base).

\subsection{Quatre centres, une même politique}

\begin{boxlecture}{Les quatre centres de LiZbiZ}
AUDIX CONSEIL cote, pour chaque centre, le contrôle clé de chacun des cinq
sous-domaines, d'après les exports de configuration remis par la DSI.

\medskip
\begin{tableaucours}{L{2.6cm}cccccc}{%
    \tableauHeader{Centre} &
    \tableauHeader{IN.1} & \tableauHeader{IN.2} & \tableauHeader{IN.3} &
    \tableauHeader{IN.4} & \tableauHeader{IN.5} & \tableauHeader{O/N/K}
}
Douala & \NON{} & \OUI{} & \OUI{} & \OUI{} & \OUI{} & 4 / 1 / 0 \\
Yaoundé & \OUI{} & \NON{} & \OUI{} & \OUI{} & \OUI{} & 4 / 1 / 0 \\
Garoua & \NON{} & \NON{} & \NON{} & \NEK{} & \NON{} & 0 / 4 / 1 \\
Bafoussam & \NON{} & \NON{} & \NON{} & \NEK{} & \NEK{} & 0 / 3 / 2 \\
Ensemble & & & & & & 8 / 9 / 3 \\
\end{tableaucours}

\medskip
\textbf{Constatation.} Ensemble~: $\mu = 0{,}40$~; $U = 0{,}15$~;
$H \approx 1{,}46$~bit. Douala et Yaoundé ont chacun un écart~; Garoua et
Bafoussam n'appliquent pratiquement pas la politique. Pour ces deux
centres, l'état des correctifs n'a pas pu être établi~: aucun export n'a
été produit.

\medskip
\textbf{Analyse} (confiance moyenne). Non-conformité claire pour
l'ensemble, mais la cause est localisée~: selon la DSI, les centres de
Garoua et de Bafoussam ont été installés par un prestataire, sans
reprendre la configuration de référence (à confirmer par les dossiers
d'installation).
Une moyenne des quatre centres serait trompeuse~; la recommandation
porte sur une configuration de référence imposée à chaque nouveau site,
et vérifiée avant sa mise en service.

\medskip
\hfill\lecture{0{,}40}{0{,}15}{1{,}46}
\end{boxlecture}

% ============================================================================
\section{Encadrer les Tests Techniques}
\label{sec:c16-tests}
% ============================================================================

\subsection{Pourquoi une autorisation écrite}

Un auditeur qui scanne un réseau ou tente une connexion accomplit
matériellement des actes que la loi réprime lorsqu'ils sont faits sans
droit.

\begin{boxjuris}
\textbf{\loicam{n°~2010/012}, article~69}~: \textit{« Est puni d'un
emprisonnement de cinq (05) à dix (10) ans et d'une amende de 10.000.000
(dix millions) à 100.000.000 (cent millions) F CFA ou de l'une de ces deux
peines seulement, celui qui accède sans droit, et en violation des mesures
de sécurité, à l'ensemble ou à une partie d'un réseau de communications
électroniques, d'un système d'information ou d'un équipement terminal,
afin d'obtenir des informations ou des données, en relation avec un
système d'information connecté à un autre système d'information. »}

\medskip
\textbf{Article~70} (résumé)~: est puni d'amende celui qui provoque par
saturation l'attaque d'une ressource de réseau ou d'un système
d'information, dans le but d'empêcher la réalisation des services
attendus.

\medskip
\textbf{Analyse.} L'autorisation écrite du responsable du système est ce
qui donne à l'auditeur le droit d'accéder. Sans elle, le test technique
peut tomber sous l'article~69. Un test de charge non encadré peut
produire l'effet visé par l'article~70~: les tests de déni de service
sont exclus, sauf autorisation expresse et environnement dédié.
\end{boxjuris}

\subsection{La note d'autorisation}

\begin{boxprocedure}
\textbf{Contenu de la note d'autorisation d'un test technique.}
\begin{enumerate}
    \item \textbf{Signataire}~: une personne habilitée à engager
          l'organisation sur les systèmes visés, et, pour un système
          hébergé, l'accord écrit de l'hébergeur.
    \item \textbf{Périmètre}~: adresses, systèmes et applications inclus~;
          systèmes explicitement exclus.
    \item \textbf{Méthodes}~: lecture de configuration, scan de ports, scan
          de vulnérabilités, test d'intrusion~; ce qui est interdit (déni
          de service, modification de données).
    \item \textbf{Fenêtres}~: dates et heures autorisées, hors des périodes
          critiques.
    \item \textbf{Contacts}~: un interlocuteur joignable pendant les tests,
          et une procédure d'arrêt immédiat.
    \item \textbf{Résultats}~: confidentialité, conservation, destruction
          à la fin de la mission (art.~14 et~61, Chapitre~\ref{chap:droit}).
\end{enumerate}
\end{boxprocedure}

\subsection{Des tests qui produisent des preuves}

Un test technique produit des sorties d'outils~: listes de ports, résultats
de scan, exports de règles. Pour être opposables, elles sont conservées
brutes, horodatées et accompagnées d'une empreinte, comme toute extraction
(Chapitre~\ref{chap:preuve}). L'interprétation (gravité, exploitabilité)
vient ensuite, dans le papier de travail.

\begin{boxalert}
\textbf{Violation typique.} Lancer un scan « rapide » du réseau de Garoua
depuis le portable de l'auditeur, la veille de la signature de la note
d'autorisation, « pour gagner du temps ».

\smallskip
\textbf{Règle~:} aucun test avant la signature, aucun test hors du
périmètre ni de la fenêtre. Un constat obtenu sans droit n'est pas
opposable, et il expose l'auditeur.
\end{boxalert}

\begin{boxmetier}{RSSI}
\textbf{Préparer un audit technique}
\begin{enumerate}
    \item Fournir la matrice de flux autorisés et la configuration de
          référence de chaque type d'équipement.
    \item Faire signer la note d'autorisation par la personne habilitée,
          avec l'accord des hébergeurs.
    \item Prévenir les équipes de surveillance, pour qu'un scan autorisé ne
          déclenche pas une réponse à incident.
    \item Produire les exports demandés, avec leur date et leur empreinte.
    \item Traiter les constats par ordre de gravité, et documenter chaque
          correction pour l'audit de suivi.
\end{enumerate}
\end{boxmetier}

\begin{boxmetier}{Auditeur SI}
\textbf{Lire un export de configuration}
\begin{enumerate}
    \item Vérifier la date, l'équipement et l'empreinte de l'export.
    \item Lister les services exposés et les comparer à la configuration de
          référence.
    \item Chercher les comptes par défaut et les protocoles en clair
          (Telnet, FTP, HTTP d'administration).
    \item Reconstituer les flux autorisés par les règles, et les comparer à
          la matrice de flux.
    \item Coter chaque contrôle par centre, avant toute agrégation.
\end{enumerate}
\end{boxmetier}

% ============================================================================

\separateur

\begin{boxresume}
\textbf{Ce qu'il faut retenir de ce chapitre~:}
\begin{itemize}
    \item \textbf{Trois audits de sécurité}~: légal (loi n°~2010/012),
          d'un SMSI (ISO/IEC~27001 et 27007), technique~; ils se
          complètent.
    \item \textbf{Le SMSI est une chaîne}~: périmètre, risques, traitement,
          mesures, évaluation~; une mesure déclarée sans preuve de
          fonctionnement est \NEK{}.
    \item \textbf{Cinq sous-domaines d'infrastructure}~: segmentation,
          durcissement, protocoles, vulnérabilités, journalisation.
    \item \textbf{La segmentation se prouve} en comparant flux autorisés et
          flux observés.
    \item \textbf{Coter par site avant d'agréger}~: une moyenne peut cacher
          deux centres hors politique.
    \item \textbf{Aucun test sans autorisation écrite}~: périmètre, méthodes,
          fenêtres, contacts, sort des résultats (art.~69 et~70).
\end{itemize}
\end{boxresume}

\bigskip

\begin{boxexo}[title={\ding{45}\quad Exercice~16.1. Quel audit~? \niveaubase}]
Pour chaque demande, indiquez l'audit de sécurité correspondant.
\begin{enumerate}[(a)]
    \item L'ANTIC annonce son passage annuel à la BANQUE SANAGA.
    \item Le RSSI veut savoir si les règles des pare-feu de Garoua
          correspondent à la matrice de flux.
    \item La BANQUE SANAGA veut un certificat pour son centre de données.
    \item Le comité d'audit demande si les écarts du dernier audit interne
          du SMSI ont été clos.
\end{enumerate}
\end{boxexo}

\begin{boxexo}[title={\ding{45}\quad Exercice~16.2. Matrice de flux \niveaumoyen}]
La politique de LiZbiZ autorise~: Internet vers DMZ en 443~; DMZ vers
applicatif en 8443~; utilisateurs vers applicatif en 8443~; administration
vers toutes les zones en 22. L'export du pare-feu de Bafoussam montre en
outre~: utilisateurs vers applicatif en 1433 et en 3389~; Internet vers
administration en 23.
\begin{enumerate}[(a)]
    \item Dressez la matrice des flux autorisés et observés.
    \item Classez les écarts par gravité, en justifiant.
    \item Quelle cote retenir pour IN.1 et pour IN.3 à Bafoussam~?
    \item Rédigez la constatation de l'observation la plus grave.
\end{enumerate}
\end{boxexo}

\begin{boxexo}[title={\ding{45}\quad Exercice~16.3. La note d'autorisation \niveauexpert}]
AUDIX CONSEIL doit réaliser un scan de vulnérabilités des quatre centres de
LiZbiZ. Le centre de Garoua est hébergé par un prestataire local.
\begin{enumerate}[(a)]
    \item Rédigez la note d'autorisation.
    \item Qui doit la signer pour Garoua, et pourquoi~?
    \item Le DSI propose d'ajouter un « test de charge » de l'application.
          Que répondez-vous, au regard de l'article~70~?
    \item Comment les résultats du scan deviennent-ils des preuves
          opposables~?
\end{enumerate}
\end{boxexo}

\bigskip

\begin{boxinfo}
\textbf{Transition.} Une infrastructure sûre doit aussi durer. Le
Chapitre~\ref{chap:audit-continuite} audite la production, la continuité
et les tiers~: exploitation, sauvegardes, restaurations, plans de
continuité et prestataires.
\end{boxinfo}
